\chapter{相关基础与研究综述}\label{ch:related}

\section{推荐系统基础}
\subsection{协同过滤与矩阵分解}
协同过滤（CF）是推荐系统的经典范式：基于"相似用户偏好相似物
品"的假设，从全体用户的交互矩阵中学习用户与物品的隐向量表示。
Sarwar 等人~\cite{sarwar2001itembased} 提出的基于物品的 CF 以
物品共现相似度为核心；Koren 等人~\cite{koren2009matrix} 的矩阵
分解（MF）将用户—物品交互矩阵分解为低维隐因子的内积。Rendle 等
人~\cite{rendle2009bpr} 的 BPR 将排序目标引入隐因子学习，以
"正样本得分高于负样本"的成对损失直接优化排序质量，成为隐式反
馈排序学习的标准基线。He 等人~\cite{he2017ncf} 进一步以多层感
知机替代内积，刻画非线性的用户—物品交互。

\subsection{序列推荐}
序列推荐将用户历史建模为按时间排序的交互序列，捕捉动态与短期
偏好。GRU4Rec~\cite{hidasi2016gru4rec} 首先将门控循环单元引入
会话推荐；SASRec~\cite{kang2018sasrec} 以因果自注意力对序列建
模，在稀疏与稠密场景均表现稳健；BERT4Rec~\cite{sun2019bert4rec}
以双向自注意力实现了更深的序列上下文。图协同方法方面，
NGCF~\cite{wang2019ngcf} 将用户—物品二部图上的高阶连通性显式编
码进协同信号；GC-MC~\cite{berg2017graphconv} 以图卷积完成矩阵
补全。

\subsection{自监督与数据增强}
自监督学习为推荐系统提供了不依赖外部标注的辅助信号。
SGL~\cite{xie2022sgl} 通过节点/边扰动构造对比视图；
SimGCL~\cite{yu2022simgcl} 证明均匀嵌入噪声即可带来与结构扰动
相当的增益；CoSeRec~\cite{liu2021coserec} 系统设计了序列级扰动
（插入/替换/截断等）并以对比目标增强序列表示。这些工作表明：
\emph{扰动感不变性}本身是推荐表示学习的重要监督来源——这一结论
是本文多视图一致性生成器的思想源头之一。

\section{递归自改进（RSI）基础}
\subsection{闭环范式}
递归自改进指模型通过自身输出持续改进自身参数的训练范式，其概
念源头可追溯至 Schmidhuber 的 Gödel 机~\cite{schmidhuber2003goedel}
与自我改进智能的讨论~\cite{yudkowsky2008rsi}。一个 RSI 系统包
含三个环节的循环：\textbf{生成}（模型产生新的训练信号）、
\textbf{评估}（系统判定信号质量）、\textbf{迭代}（以合格信号更
新参数，循环往复）。

\subsection{语言模型中的 RSI 方法谱系}
Self-Refine~\cite{madaan2023selfrefine} 让模型对自身输出给出文
字反馈并迭代精炼；STaR~\cite{zelikman2022star} 以模型自生成的
推理路径配合合理性校验构造监督信号，实现"自教推理"；
Self-Improve~\cite{huang2023selfimprove} 用自一致性多数投票产生
高置信伪标注再行微调；Reflexion~\cite{shinn2023reflexion} 将失
败轨迹转化为语言化反思存入情景记忆；
Self-Rewarding~\cite{yuan2024selfrewarding} 训练模型以 LLM-as-
Judge 的方式为自己生成奖励；SPIN~\cite{chen2024spin} 将自改进
建模为与历史版本的自博弈。这些方法的共性是：\emph{评估环节的
质量决定闭环的命运}——无质量控制的自我生成信号会引发分布收窄
与性能退化。

\section{RSI 在推荐系统中的相关尝试与局限}
将自生成信号用于推荐的研究尚处早期。可类比的三条线索：
（1）自训练式伪标注：半监督学习经典方法
~\cite{yarowsky1995wsd,lee2013pseudo} 在推荐场景的移植尝试，多
数沿用"置信度阈值"做准入，未处理推荐特有的流行度偏差；
（2）对比/扰动增强~\cite{xie2022sgl,yu2022simgcl,liu2021coserec}
：在表示空间注入不变性，但不构成跨轮次的参数自改进闭环；
（3）基于 LLM 的数据合成~\cite{wu2024llmrecsurvey}：以外部大模
型生成推荐训练数据，其生成器与被训练模型分离，不属于严格意义
的递归自改进。

\textbf{批判性总结。}现有方案的局限可归纳为三点：其一，
\emph{置信度门控的结构性偏袒}——预测置信度天然与物品流行度正
相关，纯置信度门控会系统性放行头部样本，恰是漂移之源；
其二，\emph{漂移不可见}——多数方案未定义、未度量自生成分布对
真实分布的偏离，遑论在线控制；其三，\emph{生成即自强化}——以
模型 top-$k$ 输出直接作为训练标签，生成环节本身就编码了对头部
的偏袒，门控只能在下游缓解。本文的框架正是针对这三点逐一对症
设计（分层标准化、分布预算、一致性生成）。

\section{评价指标}
\subsection{推荐准确性指标}
全量排序（full ranking）协议下，对每个用户将测试物品与全部未
交互物品比较。NDCG@$K$ 对命中位置做对数折损加权，HR@$K$ 度量
top-$K$ 命中率；二者为序列推荐的标准指标
~\cite{kang2018sasrec,sun2019bert4rec}。
\subsection{RSI 迭代质量指标}
本文额外引入刻画闭环健康的指标：\textbf{合成分布漂移}，即自生
成样本的流行度分层分布 $\bm{p}_{\mathrm{syn}}$ 对真实分布
$\bm{p}_{\mathrm{real}}$ 的 Jensen--Shannon 距离
$D_{\mathrm{JS}}$~\cite{gretton2012mmd}；\textbf{准入率}，通过
门控的候选样本比例；\textbf{迭代曲线}，指标随轮次 $r$ 的变化轨
迹，用于检测性能退化；\textbf{跨轮稳定性}，最终轮与最优轮的差
值。显著性检验采用配对 $t$ 检验。

\section{本章小结}
本章梳理了序列推荐、半监督自训练与语言模型 RSI 三条文献线，指
出推荐场景 RSI 的三点结构性局限：置信度门控的结构性偏袒、漂移
不可见、生成即自强化。第二章的批判性分析直接对应第三章框架的三
个设计决策。
